\documentclass[11pt]{article}

\usepackage[margin=1.1in]{geometry}
\usepackage{amsmath,amssymb,amsthm,mathtools,stmaryrd}
\usepackage{mathpartir}
\usepackage{tikz-cd}
\usepackage{enumitem}
\usepackage[expansion=false]{microtype}
\usepackage[hidelinks]{hyperref}

% ---- theorem environments ----
\theoremstyle{definition}
\newtheorem{definition}{Definition}[section]
\newtheorem{example}[definition]{Example}
\newtheorem{remark}[definition]{Remark}
\newtheorem{construction}[definition]{Construction}

\theoremstyle{plain}
\newtheorem{proposition}[definition]{Proposition}
\newtheorem{theorem}[definition]{Theorem}
\newtheorem{lemma}[definition]{Lemma}
\newtheorem{corollary}[definition]{Corollary}
\newtheorem{conjecture}[definition]{Conjecture}
\newtheorem{principle}[definition]{Principle}

% ---- notation ----
\newcommand{\p}{\partial}
\newcommand{\T}{T}
\newcommand{\Coll}{S}
\newcommand{\List}{\mathsf{List}}
\newcommand{\Mset}{\mathsf{Bag}}
\newcommand{\Pfin}{\mathsf{P}_{\mathrm{fin}}}
\newcommand{\Dist}{\mathsf{D}}
\newcommand{\plug}{\mathsf{plug}}
\newcommand{\Split}{\mathsf{Split}}
\newcommand{\Pack}{\mathsf{PACK}}
\newcommand{\pack}{\mathsf{pack}}
\newcommand{\Proto}{\Pi}
\newcommand{\sem}[1]{\llbracket #1 \rrbracket}
\newcommand{\real}{\Vdash}
\newcommand{\id}{\mathrm{id}}
\newcommand{\catC}{\mathcal{C}}
\newcommand{\Set}{\mathbf{Set}}
\newcommand{\op}[1]{\widehat{#1}}
\newcommand{\N}{\mathbb{N}}
\newcommand{\qmap}{q}
\newcommand{\Pos}{\mathsf{Pos}}

\title{\bfseries Packaging the Splittings\\[2pt]
\large Realizability, collection monads, and the choice content of the
McBride derivative in a finite-limits metalanguage}
\author{L.\,G.\ Meredith\\\small F1R3FLY.io}
\date{\today}

\begin{document}
\maketitle

\begin{abstract}
\noindent
i want to make visible the seams along which a logic is assembled. The
commitment that makes the seams visible is a realizability semantics:
a formula denotes the collection of computations that witness its
computational content. Under that commitment a logic factors into three
components --- a \emph{collection} component (how witnesses are aggregated),
a \emph{structural} component (term constructors lifted to predicates on
witnesses), and a \emph{behavioral} component (contexts labeling
transitions lifted to modalities). The collection component sits awkwardly
with categorical logic, because the subobject classifier bakes in
\emph{set}-shaped collection: there is no list of all lists. i reconcile the
two by going syntactic. Treating every collection as a term type, i use the
McBride derivative to form the object of all splittings $(\p\T,\T)$ --- the
\emph{proto power object} --- and then ask what must be added to the ambient
metalanguage to \emph{package} the splittings back into the collection type.
Committing the metalanguage to finite limits (cartesian logic) is what keeps
the whole apparatus syntactic: there is no subobject classifier to lean on,
so the derivative route is the only route; and choice cannot be asserted as a
proposition, so it can enter only as a named section operation. The central
claim is that packaging is exactly a section of the quotient presenting the
collection monad, that its choice content equals the assertion that the atom
object admits a linear order, and that going syntactic discharges this because
term positions are already ordered. The construction is kept parametric in the
order on atoms, isolating the choice residue at the atom boundary --- which is
where later reflective (knot-tying) constructions will act. Throughout i note
that the equations one quotients by --- commutativity and idempotence --- are
exactly the equations the no-go theorems for distributive laws turn on, so the
distributive-law obstruction and the packaging cost are one datum read twice.
\end{abstract}

\tableofcontents

\section{Why realizability, and why this matters for the argument to the
AI community}
\label{sec:motivation}

Most accounts take a logic as given: a fixed vocabulary of connectives, a
notion of model, and a truth predicate that floats above computation. Truth
values arrive from a Boolean or Heyting algebra that is postulated, not built.
For a discipline that wants its reasoning systems to be inspectable ---
traceable to constructions whose behavior can be run, audited, and held
accountable --- this is exactly the wrong starting point. The glue that binds
truth to the algebra is opaque, and once it is in place one can no longer see
what the logic is \emph{made of}.

i want the opposite. i want a logic whose every connective is an operation on
computations, so that the entire edifice is a construction and nothing is
postulated except the constructions themselves. The commitment that delivers
this is realizability. Rather than asking whether a formula is true, we ask
which computations \emph{witness} it, and we let the formula denote the
collection of its witnesses:
\[
  \sem{\varphi} \;=\; \{\, e : e \real \varphi \,\}.
\]
This is the Kleene move \cite{Kleene45,vanOosten08}, but with the aggregating
container promoted from an incidental detail to a \emph{parameter} of the
construction. Once witnesses are computations and formulae are collections of
them, every connective is forced to be an operation of one of exactly three
kinds, corresponding to the three things one can vary: how witnesses are
collected, how they are built, and how they move. This is the decomposition
i develop in \S\ref{sec:three}.

There is a second, sharper commitment, and it is the one i most want to
defend to the AI community. i insist that the ambient metalanguage --- the
language in which the collections, the constructions, and the packaging all
live --- be a \emph{finite-limits} (cartesian, essentially-algebraic) theory
\cite{Freyd72,GabrielUlmer71,PalmgrenVickers07}. The reason is not
minimalism for its own sake. It is that a finite-limits metalanguage is the
maximal commitment to ``everything is a construction'':
\begin{itemize}[leftmargin=1.4em]
\item there are no truth values floating in from a Boolean algebra, because
      there is in general no subobject classifier (a lex category need not be
      a topos);
\item there is no non-constructive existential, because the only existentials
      one may assert are provably \emph{unique} ones, carved out as partial
      operations on equationally-defined subdomains;
\item and --- the point that will do the real work in
      \S\ref{sec:finlim} and \S\ref{sec:packaging} --- choice cannot be
      \emph{stated} as a proposition at all. There is no $\forall$, and $\exists$
      is available only when the witness is unique. So the only way to
      ``add choice'' is to add a \emph{function symbol} to the signature,
      together with equations making it a section. In a finite-limits
      metalanguage, choice is not a truth one assumes; it is an operation one
      names.
\end{itemize}
This is exactly the inspectability i am after. When choice enters, it enters
as a symbol you can point at, with a documented arity and documented
equations, and you can localize precisely where in the construction it was
spent. Nothing is hidden in a semantic ether. The seams are visible because
we have refused the glue --- the subobject classifier --- that ordinarily
hides them.

The remainder of the note carries this out. \S\ref{sec:three} gives the
three-component decomposition. \S\ref{sec:dial} shows that the collection
component is a substructural dial whose settings are governed by two
equations, commutativity and idempotence. \S\ref{sec:disconnect} states the
disconnect with categorical logic, and \S\ref{sec:nogo} shows that the same
two equations are what the no-go theorems for distributive laws turn on ---
so the obstruction to a uniform collection component and the cost of packaging
are, as we will see, one datum read twice. \S\ref{sec:finlim} cashes out the
finite-limits commitment. \S\ref{sec:derivative} builds the proto power
object from the McBride derivative, and \S\ref{sec:packaging} defines
packaging as a section and states the choice conjecture, parametric in the
order on atoms. \S\ref{sec:grading} places the conjecture in the Weihrauch
lattice, \S\ref{sec:instances} works the three instances, and
\S\ref{sec:discussion} draws out the consequences, including the absence of
any classicality collapse and the reflective use of the atom-order parameter.

\section{Three components of a logic under realizability}
\label{sec:three}

Fix a representation of computation. Throughout i take it to be a term
rewrite system: terms are computational states, and rewrites are steps. Write
$\T$ for the type of terms. Fix also a way of aggregating witnesses: a
\emph{collection type} $\Coll$, which will always carry the structure of a
monad (\S\ref{sec:dial}). A formula denotes an object $\sem{\varphi} :
\Coll(\T)$ --- a collection of terms.

Under these two fixings, three families of connectives arise, and every
connective is of one of the three kinds.

\paragraph{Collection connectives.}
Any operation on $\Coll$ lifts to a connective. If witnesses are collected in
sets, the semilattice operations of $\Coll = \Pfin$ give the familiar
connectives:
\[
  \sem{\varphi \wedge \psi} = \sem{\varphi} \cap \sem{\psi},
  \qquad
  \sem{\varphi \vee \psi} = \sem{\varphi} \cup \sem{\psi}.
\]
This is where the choice of container is felt: change $\Coll$ from sets to
multisets, lists, trees, or distributions, and these connectives change with
it. The collection component is the parameter i most want to vary, and
\S\ref{sec:dial} explains what varying it does.

\paragraph{Structural connectives.}
Each term constructor of $\T$ lifts to a predicate transformer. If
$f : \T^{n} \to \T$ is an $n$-ary constructor, define
\[
  \sem{\op{f}(\varphi_{1},\dots,\varphi_{n})}
  \;=\;
  \{\, f(e_{1},\dots,e_{n}) : e_{i} \real \varphi_{i} \,\}.
\]
The lifted constructor $\op{f}$ is a \emph{structural} predicate: it speaks
about the \emph{shape} of a witness, asserting that the witnessing computation
was built by $f$ from sub-witnesses of the $\varphi_{i}$. In the reflective
setting these are the constructors of the rho calculus lifted to namespace
predicates \cite{MeredithRadestock05a,MeredithRadestock05b}; in general they
are the introduction half of an Operational-Semantics-in-Logical-Form (OSLF)
presentation.

\paragraph{Behavioral connectives.}
Each rewrite is labeled by the context in which it fires. A step
$C[l] \to C[r]$ carries the one-hole context $C$, an element of the
\emph{derivative} $\p\T$ (\S\ref{sec:derivative}). Lifting the labels to
modalities gives
\[
  \sem{\langle C\rangle \varphi}
  = \{\, e : e \xrightarrow{\,C\,} e',\ e' \real \varphi \,\},
  \qquad
  \sem{[C]\varphi}
  = \{\, e : \forall e'.\ e \xrightarrow{\,C\,} e' \Rightarrow e' \real \varphi \,\},
\]
the context-indexed Hennessy--Milner modalities. This is the elimination half
of OSLF.

\begin{remark}
The structural and behavioral components are the two halves of OSLF, generated
respectively from constructors and from labeled transitions. They are already
well-understood in the reflective setting. The component that resists the
categorical-logic packaging, and that this note is about, is the first one ---
collection --- together with the derivative machinery ($\p\T$) that turns out
to glue it to the third.
\end{remark}

\section{The collection component is a substructural dial}
\label{sec:dial}

The collection type $\Coll$ is a monad, and the algebraic theory presenting
$\Coll$ is not incidental: \emph{its equations are the structural rules of the
propositional logic the collection component induces.} This is the first fact
that makes the decomposition worth having, because it says the collection
component is not ``logic plus a data-structure flavor''; it \emph{is} the
substructural hierarchy, indexed by how far $\Coll$ sits from set-shaped.

Write the binary aggregation of $\Coll$ as $\ast$ with unit $\varepsilon$.
The correspondence is:
\[
\begin{array}{lll}
  \text{commutativity of } \ast & \longleftrightarrow & \text{exchange},\\[2pt]
  \text{idempotence of } \ast    & \longleftrightarrow & \text{contraction},\\[2pt]
  \text{unit/absorption laws}    & \longleftrightarrow & \text{weakening}.
\end{array}
\]
Reading the standard collection monads through this correspondence:
\begin{itemize}[leftmargin=1.4em]
\item $\Pfin$, the finite-powerset monad --- a commutative, idempotent monoid
      (a join-semilattice) --- gives exchange and contraction (and, with the
      unit, weakening): the full structural regime, hence Heyting/Boolean
      propositional logic.
\item $\Mset$, the finite-multiset monad --- the free commutative monoid,
      commutative but not idempotent --- gives exchange without contraction:
      the multiplicative-linear regime with $\otimes$.
\item $\List$, the list monad --- the free monoid, neither commutative nor
      idempotent --- gives neither exchange nor contraction: a
      non-commutative, Lambek-style regime.
\item $\Dist$, the distribution monad --- commutative with a convex (barycentric)
      structure in place of idempotence --- gives the probabilistic regime.
\end{itemize}

\begin{principle}[Collection is substructure]
\label{prin:dial}
The setting of the collection dial is the pair of Booleans
$(\text{commutative?}, \text{idempotent?})$ of $\Coll$, and this pair is
exactly the presence or absence of exchange and contraction in the induced
logic.
\end{principle}

Keep the two coordinates of Principle~\ref{prin:dial} in view. They are the
same two coordinates on which the no-go theorems of \S\ref{sec:nogo} turn,
and the same two equations one quotients by when packaging in
\S\ref{sec:packaging}. Everything in this note is, at bottom, bookkeeping on
commutativity and idempotence.

\section{The disconnect with categorical logic}
\label{sec:disconnect}

Categorical logic packages the collection component through the subobject
classifier $\Omega$: the power object of $A$ is $\Omega^{A}$, the object of
subobjects of $A$. This is a beautiful and total solution --- for one setting
of the dial. It hard-wires $\Coll = $ powerset. Two things go wrong when we
try to honor Principle~\ref{prin:dial} and turn the dial.

First, the classifier fixes the collection shape. There is no
$\Omega$-parametrized-by-$\Coll$; the power object is subsets, and only
subsets. One cannot, inside the $\Omega$ story, ask for the multiset or list
of witnesses.

Second, and more tellingly, even the operation of forming ``the collection of
all subcollections \emph{as an object of the same shape}'' fails away from
sets. The slogan is: \emph{there is no list of all lists.} To assemble all
sublists of a given list into a single \emph{list}, one must linearly order
the sublists; there is no canonical order on them, and choosing one is choice.
For sets the analogous assembly is free of order --- the subsets form a set
with nothing to arrange --- which is exactly why the $\Omega$ story works for
the powerset setting and only there.

\begin{remark}
The disconnect is not a defect of categorical logic to be repaired by a
cleverer classifier. It is a signal that the packaging of the collection
component into an object of the same shape is \emph{order-laden}, and that the
$\Omega$ story succeeds precisely by living at the one setting of the dial
(powerset) where the order requirement is vacuous. The reconciliation in
\S\ref{sec:derivative}--\S\ref{sec:packaging} takes the order requirement
seriously and localizes it, rather than wishing it away.
\end{remark}

\section{The same two equations: distributive laws and their no-go theorems}
\label{sec:nogo}

Before the reconciliation, i want to record why the collection dial cannot be
turned freely if one insists on the standard monadic route to the structural
connectives, because it is the same two equations again.

The structural connectives of \S\ref{sec:three}, viewed monadically, ask the
collection monad $\Coll$ to commute past the term monad $\T$. Lifting a
constructor $f : \T^{n} \to \T$ to collections requires a coherent map
$(\Coll\,A)^{n} \to \Coll(A^{n}) \to \Coll\,A$, and for this to be coherent
across \emph{all} the constructors of $\T$ simultaneously one wants a
distributive law
\[
  d : \Coll \circ \T \;\Rightarrow\; \T \circ \Coll,
\]
in the sense of Beck \cite{Beck69}. A distributive law is exactly what lets
$\Coll$-collections be pushed through $\T$-structure functorially, yielding
the lifted constructors.

The trouble is that such laws are scarce, and the scarcity is governed by
commutativity and idempotence. The covariant powerset monad does not
distribute over itself \cite{KlinSalamanca18}; the distribution monad does not
distribute over powerset \cite{VaraccaWinskel06}; and the general no-go
theorems of Zwart and Marsden \cite{ZwartMarsden19} show that once the two
theories being combined carry equational shapes of the wrong kind --- an
idempotence or an absorption clashing with the other theory --- \emph{no}
distributive law can exist. The upshot is that, along the distributive-law
route, the only collection monads that commute nicely past term-like $\T$ are
the set/multiset/distribution-shaped ones: exactly the ones near the
top-and-commutative corner of the dial.

\begin{principle}[One datum, read twice]
\label{prin:twice}
The equations that obstruct the distributive law $d : \Coll\T \Rightarrow
\T\Coll$ --- commutativity and idempotence --- are the same equations that,
in \S\ref{sec:packaging}, one must quotient $\List$ by to present $\Coll$, and
hence the same equations whose section must be postulated to package. The
no-go obstruction and the packaging cost are one datum read twice: once as a
law that fails, once as an operation that must be added.
\end{principle}

The reconciliation of \S\ref{sec:derivative}--\S\ref{sec:packaging} is, in
effect, a third way around the no-go theorems. The two known escapes are to
\emph{quotient} (convex powerset, indexed valuations --- adjust $\Coll$ until
the equations agree) or to \emph{restrict} (finite, ordered --- shrink until
the clash cannot arise). The route here is to \emph{add a global structuring
operation} --- the section --- and pay for it in a named, localized axiom.
\S\ref{sec:discussion} argues these are dual dodges.

\section{The finite-limits metalanguage does exactly what we need}
\label{sec:finlim}

Now the metalanguage commitment earns its keep. Take the ambient metalanguage
to be a finite-limits theory: the cartesian fragment, with $\top$, $\wedge$,
$=$, and existentials only where the witness is provably unique
\cite{Freyd72,GabrielUlmer71,PalmgrenVickers07}. Absent are $\vee$, $\bot$,
$\Rightarrow$, $\neg$, unrestricted $\exists$ and $\forall$, and --- in
general --- a subobject classifier. Three consequences, each of which turns a
would-be obstruction into a feature.

\paragraph{(1) There is no $\Omega$ to lean on, so the syntactic route is the
only route.} In a topos one might hope to repair \S\ref{sec:disconnect} with a
cleverer use of the classifier. In a finite-limits metalanguage that hope is
simply unavailable: a lex category need not be a topos, and the power object
is not on the table. The McBride-derivative construction of
\S\ref{sec:derivative} is therefore not \emph{an} alternative to the
$\Omega$-story; it is the \emph{only} story. The commitment vindicates going
syntactic instead of competing with it.

\paragraph{(2) Choice cannot be a proposition, only an operation.} In
cartesian logic one cannot write the choice principle in $\forall\exists$
form: there is no $\forall$, and $\exists$ may be asserted only under
uniqueness. So ``adding choice'' cannot mean asserting a sentence. It can only
mean enlarging the signature with a function symbol and equations that make it
a section of some quotient. This is precisely the shape that packaging will
take in \S\ref{sec:packaging}: $\pack$ is a new operation, and the equation
$\qmap \circ \pack = \id$ is choice-as-structure. The metalanguage forces the
conjecture into its most honest and most inspectable form.

\paragraph{(3) No classicality collapse.} In a topos, internal choice forces
Boolean logic (Diaconescu \cite{Diaconescu75,MacLaneMoerdijk92}). That
argument needs $\vee$, comprehension via the power object, and effective
quotients to run --- \emph{none} of which are present in the cartesian
fragment. Consequently one may add packaging/choice as pure structure without
collapsing the object logic to classical. Choosing finite limits is exactly
what keeps choice \emph{structural} rather than \emph{logical}. This is the
technical heart of the inspectability argument of \S\ref{sec:motivation}: the
choice we add is a nameable operation and it does not silently classicalize
everything downstream.

\section{The McBride derivative and the proto power object}
\label{sec:derivative}

i now build the object that replaces the power object. The move is to treat
every collection \emph{syntactically} --- as a term type --- and to use the
McBride derivative to enumerate its splittings.

\subsection{Terms, derivatives, and the zipper}

Let $F$ be a polynomial (container) endofunctor and let $\T = \mu X.\,F X$ be
its initial algebra, the type of terms. Write $\p F$ for the derivative of
$F$: the functor of one-hole $F$-contexts \cite{McBride01,AAGM05,Huet97}. The
derivative of the fixpoint is a list of layers up the spine:
\[
  \p\T \;\cong\; \bigl(\p F[\T]\bigr)^{\!\ast}
       \;=\; \List\!\bigl(\p F[\T]\bigr),
\]
each layer an $F$-node with one hole and its other children filled by
subtrees. Plugging fills the hole:
\[
  \plug : \p\T \times \T \longrightarrow \T,
  \qquad
  \plug(\mathsf{nil}, t) = t, \quad
  \plug(\ell :: c, t) = \plug\bigl(c, \mathsf{fill}(\ell, t)\bigr).
\]

\begin{definition}[Splittings]
The object of \emph{splittings} of $\T$ is
\[
  \Split(\T) \;=\; \{\, (c,t) \in \p\T \times \T \,\}
             \;\cong\; \textstyle\sum_{c:\p\T} \T,
\]
the pointed terms, with the whole recovered by $\plug(c,t)$. A splitting is a
context together with a focus.
\end{definition}

Two facts about splittings are load-bearing.

\begin{proposition}[Positions are ordered by syntax]
\label{prop:dewey}
The positions of a term --- equivalently, the contexts $c$ appearing in its
splittings --- carry a canonical linear order, the traversal (Dewey) order,
definable from the term's own structure by a fold. Hence $\p\T$ is linearly
ordered \emph{per term}, with no appeal to any order on the atoms.
\end{proposition}

\begin{proof}[Sketch]
A depth-first left-to-right traversal assigns to each position a Dewey path,
and Dewey paths are linearly ordered lexicographically. The assignment is a
catamorphism over $\T$, so it is definable once $\T$ is equipped with its
recursor (\S\ref{sec:packaging}); no comparison on atoms is used, because
positions are indexed by constructor-argument slots, which are ordered by the
arity, not by the atoms sitting in them.
\end{proof}

\begin{proposition}[The aperture is the modal label]
\label{prop:glue}
The contexts that index splittings are the same objects that label rewrites in
the behavioral component of \S\ref{sec:three}: both are elements of $\p\T$.
Thus the collection component (via its splittings) and the behavioral
component (via its modal labels) are glued at the derivative. In the agency
vocabulary, the aperture --- the locus at which a computation may be acted
upon --- is simultaneously the modal label and the raw material of the proto
power object.
\end{proposition}

Proposition~\ref{prop:glue} is the structural dividend of going through $\p\T$:
components (1) and (3) of \S\ref{sec:three} are not merely analogous, they
share their carrier.

\subsection{Collections are term types too}

The syntactic stance says: present the collection type itself as a term type,
so the same derivative machinery applies to it. The paradigm is the list
monad,
\[
  \List\,A = \mu X.\, 1 + A \times X,
\]
a genuine $\mu$-type. Its derivative in the element parameter is the pair of a
prefix and a suffix,
\[
  \p_{A}(\List\,A) \;\cong\; \List\,A \times \List\,A,
\]
so a splitting of a list is a triple $(\text{prefix}, \text{focus},
\text{suffix})$ --- exactly a choice of one element to keep in view, with its
surroundings. Gathering foci across a family of splittings and re-collecting
rebuilds a \emph{sub}collection. This is the concrete content of ``no list of
all lists'': each individual sublist is free to form, but assembling all of
them into one list requires ordering them.

\begin{definition}[Proto power object]
\label{def:proto}
For a term type $\T$, the \emph{proto power object} $\Proto(\T)$ is the
container of splittings, written $(\p\T, \T)$: shapes are contexts $c : \p\T$,
and the position at each shape is the focus of type $\T$. Its extension is
$\Proto(\T)(X) = \sum_{c:\p\T} X$, and $\Proto(\T)(\T) = \Split(\T)$.
\end{definition}

$\Proto(\T)$ is ``proto'' because it enumerates one-step focusings; it is not
yet closed under gathering foci into a subcollection of the same shape as
$\T$. Closing it is packaging.

\section{Packaging as a section, and the choice conjecture}
\label{sec:packaging}

\subsection{Presenting collections as quotients of the free case}

Present each collection monad $\Coll$ syntactically, as a quotient of the
free (list) collection:
\[
  \List\,A \xrightarrow{\ \qmap_{\Coll}\ } \Coll\,A,
\]
where
\[
  \qmap_{\List} = \id, \qquad
  \qmap_{\Mset} = (-)/\text{perm}, \qquad
  \qmap_{\Pfin} = (-)/(\text{perm} + \text{idem}).
\]
$\List$ is the free monoid --- the syntactic collection, paying nothing.
$\Mset$ adds commutativity (permutation). $\Pfin$ adds commutativity and
idempotence. These are precisely the two coordinates of the dial
(Principle~\ref{prin:dial}) and the two culprits of the no-go theorems
(Principle~\ref{prin:twice}).

\subsection{Two additions before choice}

Two things must be added to the bare finite-limits base before packaging can
even be discussed, and honesty requires naming them.

\begin{remark}[The recursor is the first non-Horn addition]
\label{rem:recursor}
Having ``the type of terms'' means an initial algebra $\mu F$, and using it
--- defining any fold, hence $\pack$, hence the Dewey order of
Proposition~\ref{prop:dewey} --- requires its \emph{recursor}. Formation and
introduction for a $\mu$-type are algebraic; elimination (induction/recursion)
is not cartesian. So the recursor is the first thing added to the bare base,
and we add the \emph{weakest} one that still defines $\pack$: a
non-dependent catamorphism. Everything downstream, including packaging, is a
fold and inherits exactly this recursor and no more.
\end{remark}

\begin{remark}[Quotients before sections]
For non-free $\Coll$, the map $\qmap_{\Coll}$ is a coequalizer, which a bare
lex category may lack. Presenting $\Coll$ therefore already commits one to
effective quotients (pretopos-flavored structure) \emph{before} asking them to
split. Packaging then splits what one was forced to construct.
\end{remark}

So the layering is: cartesian base $\;\to\;$ weakest recursor for $\T$
(Remark~\ref{rem:recursor}) $\;\to\;$ effective quotients for non-free $\Coll$
$\;\to\;$ the section below. Choice sits only at the top layer.

\subsection{The packaging axiom}

\begin{definition}[Packaging]
\label{def:pack}
A \emph{packaging} for $\Coll$ is a signature operation
\[
  \pack_{\Coll} : \Coll\,A \longrightarrow \List\,A
  \qquad\text{with}\qquad
  \qmap_{\Coll} \circ \pack_{\Coll} = \id_{\Coll A},
\]
a section of the presenting quotient: a choice, for each $\Coll$-collection, of
a list representative. The axiom asserting the existence of $\pack_{\Coll}$ is
$\Pack_{\Coll}$. In the finite-limits metalanguage $\Pack_{\Coll}$ is not a
proposition but the enlargement of the signature by $\pack_{\Coll}$ and the
above equation --- choice-as-structure, per \S\ref{sec:finlim}(2).
\end{definition}

Definition~\ref{def:pack} is the syntactic form of the closure operation on
the proto power object: to assemble the splittings of $\Proto(\T)$
(Definition~\ref{def:proto}) into a $\Coll$-collection of the same shape is to
apply $\qmap_{\Coll}$, and to make that assembly canonical --- to build
$\T[(\p\T,\T)]$, the collection whose atoms are splittings --- is to choose the
section $\pack_{\Coll}$.

\subsection{The choice conjecture, parametric in the atom order}

\begin{theorem}[Derivability of packaging, per instance]
\label{thm:instances}
Modulo the weakest recursor of Remark~\ref{rem:recursor}, and for the section
taken uniformly (natural in $A$):
\begin{enumerate}[label=\textup{(\roman*)},leftmargin=2.2em]
\item $\Pack_{\List}$ is derivable outright; $\pack_{\List} = \id$.
\item $\Pack_{\Mset}$ is derivable if the atom object $A$ carries a linear
      order $\le$ definable in the theory --- take $\pack_{\Mset} =
      \mathsf{sort}_{\le}$ --- and, conversely, a uniform section of
      $\qmap_{\Mset}$ determines a linear order on $A$ (read off from its
      action on two-element bags). So $\Pack_{\Mset}$ is inter-derivable with
      ``$A$ is linearly ordered.''
\item $\Pack_{\Pfin}$ is derivable if $A$ is linearly ordered \emph{and}
      equality on $A$ is available to deduplicate --- take
      $\pack_{\Pfin} = \mathsf{dedup} \circ \mathsf{sort}_{\le}$; conversely a
      uniform section determines both an order and a canonical
      representative.
\end{enumerate}
\end{theorem}

\begin{proof}[Sketch]
(i) is immediate. For (ii), $\mathsf{sort}_{\le}$ is an insertion sort, a fold
over $\List\,A$ using $\le$, hence definable from the weakest recursor and
$\le$; it lands in a canonical (sorted) representative of each permutation
class, so $\qmap_{\Mset}\circ\mathsf{sort}_{\le} = \id$. Conversely, given a
uniform section $s$, its value on the bag $\{a,b\}$ is one of the lists
$[a,b]$ or $[b,a]$; uniformity (naturality in $A$) forces this decision to
cohere into a total order, and the monoid equations force transitivity.
(iii) adds idempotence: a section of the combined quotient must additionally
select, from each finite subset, a duplicate-free list, i.e.\ a decision of
equality plus the order of (ii). \end{proof}

Theorem~\ref{thm:instances} localizes the choice content exactly. Read
against the earlier, topos-flavored heuristic in which ordered collections
were the expensive case (linearize an unordered carrier $=$ well-order $=$
full choice), it \emph{inverts} the cost: here $\List$ is free and the
quotients are dear. The two framings concern different maps, and the
finite-limits, syntactic stance is what reconciles them:

\begin{principle}[The McBride move trades and discharges]
\label{prin:trade}
Going through the derivative trades the \emph{well-ordering flavor} of choice
for the \emph{quotient-section flavor}, and then discharges the former for
free. The foci of a term-splitting are \emph{positions}, and positions are
pre-ordered by the term's own structure (Proposition~\ref{prop:dewey}). Syntax
hands you the linear order. So over a term rewrite system, packaging is
derivable from the positional order, and choice re-enters only through the
quotient equations, and only when the atom object $A$ carries no
cartesian-definable linear order.
\end{principle}

\begin{conjecture}[Choice content of packaging]
\label{conj:main}
The choice content of $\Pack_{\Coll}$ equals the assertion that the atom
object $A$ admits a linear order. Equivalently: with an order on $A$ supplied,
$\Pack_{\Coll}$ is derivable for every finitary $\Coll$ presented as a quotient
of $\List$; without one, $\Pack_{\Coll}$ is the added choice operation, and its
strength is graded by the equations $\qmap_{\Coll}$ quotients $\List$ by. In
particular the construction is kept \emph{parametric in an order on atoms}:
supplying the parameter makes packaging free; withholding it isolates the
entire choice residue at the atom object.
\end{conjecture}

The parametricity is deliberate and is the hook for later work
(\S\ref{sec:discussion}): in a reflective setting the atoms are quoted
processes --- names --- and an order on names is exactly a place where a knot
may be tied. Leaving the order parametric localizes the choice content at the
reflection boundary, where we will want to act on it.

\section{Grading the conjecture in the Weihrauch lattice}
\label{sec:grading}

Conjecture~\ref{conj:main} is not a switch but a dial, and the natural home
for its grading is the Weihrauch lattice \cite{BrattkaGherardi11}, which
measures the uniform strength of a choice principle. i state the intended
grading as a program; the finitary cases are within reach of the sketch of
Theorem~\ref{thm:instances}, the infinitary ones are open.

\begin{conjecture}[Weihrauch grading]
\label{conj:weihrauch}
The Weihrauch degree of $\Pack_{\Coll}$ over an atom object $A$ is monotone in
the substructural strength of $\Coll$ --- i.e.\ in the equations
$\qmap_{\Coll}$ adds to $\List$ --- and is bounded as follows:
\begin{itemize}[leftmargin=1.4em]
\item $\Pack_{\List}$: the identity; the bottom of the lattice.
\item $\Pack_{\Mset}$ over finite bags on discrete $A$: the finite/bounded
      choice needed to order each finite support uniformly; low in the
      lattice, below $\mathsf{LPO}$ when $A$ is discrete.
\item $\Pack_{\Mset}$ over an arbitrary (possibly infinite) atom object: a
      closed-choice principle $\mathsf{C}_{A}$ on the carrier, climbing toward
      well-ordering as $A$ grows; the well-ordering flavor of
      Principle~\ref{prin:trade} re-emerges exactly here, when the syntactic
      order is unavailable.
\item $\Pack_{\Pfin}$: the degree of $\Pack_{\Mset}$ joined with the
      equality/deduplication decision, adding an $\mathsf{LPO}$-flavored
      component for the idempotence quotient.
\end{itemize}
\end{conjecture}

The value of the Weihrauch reading is that it reframes the no-go theorems of
\S\ref{sec:nogo} not as yes/no obstructions but as \emph{choice-strength
measurements}. Where the classical no-go statement says ``no distributive law
exists,'' the graded statement says ``the section that stands in for the
missing law has degree $d$,'' and $d$ is a monotone invariant of how far
$\Coll$ sits from the free case. That is the stronger thesis hiding inside the
bare conjecture.

\section{The three instances, concretely}
\label{sec:instances}

i collect the worked instances. In each, the theory is the cartesian base
extended by the weakest recursor for the relevant $\mu$-type
(Remark~\ref{rem:recursor}); the atom order is the parameter of
Conjecture~\ref{conj:main}.

\begin{center}
\renewcommand{\arraystretch}{1.35}
\begin{tabular}{@{}llll@{}}
\hline
$\Coll$ & equations added to $\List$ & $\pack_{\Coll}$ & choice residue \\
\hline
$\List$  & none                       & $\id$                                   & none \\
$\Mset$  & commutativity              & $\mathsf{sort}_{\le}$                   & order on $A$ \\
$\Pfin$  & commutativity, idempotence & $\mathsf{dedup}\circ\mathsf{sort}_{\le}$ & order on $A$, $=$ on $A$ \\
\hline
\end{tabular}
\end{center}

\begin{example}[List: nothing is spent]
$\List\,A = \mu X.\,1 + A\times X$; $\qmap_{\List}=\id$; the section is the
identity. The proto power object $\Proto(\List\,A)$ has $\p(\List\,A) \cong
\List\,A \times \List\,A$ (prefix, suffix), and a splitting is
$(\text{prefix},\text{focus},\text{suffix})$. Packaging is free because there
is nothing to reassemble past the given list order.
\end{example}

\begin{example}[Bag: order on atoms is spent]
$\Mset\,A = \List\,A/\text{perm}$. With $\le$ on $A$, insertion sort is a fold
selecting the sorted representative, and $\qmap_{\Mset}\circ\mathsf{sort}_{\le}
= \id$. Withhold $\le$ and no cartesian term produces a representative
uniformly; $\pack_{\Mset}$ must be postulated, and by Theorem~\ref{thm:instances}(ii)
its postulation is inter-derivable with linearly ordering $A$. This is the
cleanest witness for Conjecture~\ref{conj:main}.
\end{example}

\begin{example}[Finite powerset: order and equality]
$\Pfin\,A = \List\,A/(\text{perm}+\text{idem})$. Packaging needs the sorted
representative of the previous example \emph{and} a deduplication, hence
decidable equality on $A$ in addition to the order. The two additions are
exactly the two quotient equations, confirming Principle~\ref{prin:twice}: the
idempotence that obstructs $\Pfin$'s distributive laws is the same idempotence
that must be paid for, as $\mathsf{dedup}$, at packaging time.
\end{example}

\section{Discussion}
\label{sec:discussion}

\paragraph{No classicality collapse.} The absence of Diaconescu's argument in
the cartesian fragment (\S\ref{sec:finlim}(3)) is what makes this entire
program viable. We add a choice operation --- $\pack_{\Coll}$ --- and the
object logic does not thereby go Boolean, because the collapse needs $\vee$,
comprehension, and effective-quotient machinery that the metalanguage does not
provide. Choice remains structural. A logic assembled this way can have any
setting of the substructural dial (Principle~\ref{prin:dial}) and still admit
packaging, without being forced to classicalize. This is the payoff i most
want the AI community to register: inspectable choice that does not
contaminate the logic it serves.

\paragraph{Quotient versus choice as dual dodges.} \S\ref{sec:nogo} noted two
standard escapes from the no-go theorems --- quotient (adjust $\Coll$ until the
equations agree) and restrict (shrink until they cannot clash) --- and offered
a third, adding a global section. i conjecture these are dual: the
quotient-dodge builds a coequalizer, the section-dodge splits it; the
restrict-dodge builds a well-founded (finite, ordered) carrier, the
section-dodge well-orders it. In slogan form: \emph{coequalizer versus Zorn}.
Making this duality precise --- exhibiting the section-dodge as the formal dual
of the quotient-dodge in the appropriate fibration --- is the natural next
technical step, and it would situate the construction cleanly against the
existing distributive-law literature.

\paragraph{Graphs are the genuine frontier.} Everything here rests on $\T$
being a polynomial/container functor, so that $\p\T$ exists and
Definition~\ref{def:proto} is available. Lists and trees are fine. Graphs are
not free/regular in the naive sense, the McBride derivative does not directly
apply, and the proto power object is not defined. i bracket the graph case
explicitly rather than let it ride along with lists and trees; adapting the
derivative to non-regular carriers (or replacing it by a suitable
comonadic/traced structure) is the honest open problem for extending the
collection dial past the container-shaped settings.

\paragraph{The atom-order parameter and knot-tying.} Conjecture~\ref{conj:main}
keeps the order on atoms parametric on purpose. In the reflective (rho) setting
the atoms are quoted processes --- names --- and the boundary between process
and name is exactly the reflection boundary. An order on names is therefore a
structure imposed \emph{at} that boundary, and it is precisely the datum a
knot-tying construction will manipulate. By isolating the entire choice residue
of packaging in the atom order, the present construction leaves the knot
exactly one place to be tied: supply the order (and packaging trivializes; the
knot is untied) or withhold it (and the choice residue is the knot). This is
the sense in which the finite-limits, derivative-based account is not merely a
reconciliation of the collection component with categorical logic, but a setup
for the reflective constructions to come.

\paragraph{Summary.} The three-component decomposition makes the seams of a
logic visible under realizability; the collection component is a substructural
dial governed by commutativity and idempotence; those same two equations
obstruct the distributive laws and set the cost of packaging; the
finite-limits metalanguage removes the subobject classifier (forcing the
syntactic route) and forbids stating choice as a proposition (forcing it to be
a named section), while blocking the classicality collapse that would
otherwise make this untenable; the McBride derivative supplies the proto power
object and, through the positional order on splittings, discharges the
well-ordering flavor of choice; and what remains --- the choice content of
packaging --- is exactly the order on atoms, kept parametric for the knot to
act on. That is the construction i wanted to formalize, and the shape of the
program it opens.

\begin{thebibliography}{99}

\bibitem{Kleene45} S.~C.\ Kleene. On the interpretation of intuitionistic
number theory. \emph{Journal of Symbolic Logic}, 10:109--124, 1945.

\bibitem{vanOosten08} J.\ van Oosten. \emph{Realizability: An Introduction to
its Categorical Side}. Studies in Logic and the Foundations of Mathematics
152, Elsevier, 2008.

\bibitem{MeredithRadestock05a} L.~G.\ Meredith and M.\ Radestock. A reflective
higher-order calculus. \emph{Electronic Notes in Theoretical Computer
Science}, 141(5):49--67, 2005.

\bibitem{MeredithRadestock05b} L.~G.\ Meredith and M.\ Radestock. Namespace
logic: A logic for a reflective higher-order calculus. In \emph{Trustworthy
Global Computing (TGC)}, LNCS 3705, 2005.

\bibitem{Beck69} J.\ Beck. Distributive laws. In \emph{Seminar on Triples and
Categorical Homology Theory}, LNM 80, pp.\ 119--140, Springer, 1969.

\bibitem{KlinSalamanca18} B.\ Klin and J.\ Salamanca. Iterated covariant
powerset is not a monad. In \emph{Mathematical Foundations of Programming
Semantics (MFPS)}, ENTCS 341, 2018.

\bibitem{VaraccaWinskel06} D.\ Varacca and G.\ Winskel. Distributing
probability over non-determinism. \emph{Mathematical Structures in Computer
Science}, 16(1):87--113, 2006.

\bibitem{ZwartMarsden19} M.\ Zwart and D.\ Marsden. No-go theorems for
distributive laws. In \emph{Logic in Computer Science (LICS)}, 2019.
(Extended version: \emph{Logical Methods in Computer Science}, 2022.)

\bibitem{McBride01} C.\ McBride. The derivative of a regular type is its type
of one-hole contexts. Unpublished manuscript, 2001.

\bibitem{AAGM05} M.\ Abbott, T.\ Altenkirch, N.\ Ghani, and C.\ McBride.
$\partial$ for data: Differentiating data structures. \emph{Fundamenta
Informaticae}, 65(1--2):1--28, 2005.

\bibitem{Huet97} G.\ Huet. The zipper. \emph{Journal of Functional
Programming}, 7(5):549--554, 1997.

\bibitem{Freyd72} P.\ Freyd. Aspects of topoi. \emph{Bulletin of the
Australian Mathematical Society}, 7:1--76, 1972.

\bibitem{GabrielUlmer71} P.\ Gabriel and F.\ Ulmer. \emph{Lokal
pr\"asentierbare Kategorien}. LNM 221, Springer, 1971.

\bibitem{PalmgrenVickers07} E.\ Palmgren and S.~J.\ Vickers. Partial Horn
logic and cartesian categories. \emph{Annals of Pure and Applied Logic},
145(3):314--353, 2007.

\bibitem{Diaconescu75} R.\ Diaconescu. Axiom of choice and complementation.
\emph{Proceedings of the American Mathematical Society}, 51:176--178, 1975.

\bibitem{MacLaneMoerdijk92} S.\ Mac Lane and I.\ Moerdijk. \emph{Sheaves in
Geometry and Logic: A First Introduction to Topos Theory}. Springer, 1992.

\bibitem{BrattkaGherardi11} V.\ Brattka and G.\ Gherardi. Weihrauch degrees,
omniscience principles and weak computability. \emph{Journal of Symbolic
Logic}, 76(1):143--176, 2011.

\end{thebibliography}

\end{document}
